\section{Transport by capped chains}\label{sew:section}

We now combine the local calculation of \cref{fp:local-algebra}, the
evaluation correspondences of \cref{ev:restriction,ev:groups}, and the cap
identity of \cref{cap:identity}.  The conclusion concerns tests of an
absolute Virasoro error.  The relative theories entering its proof are
ordinary relative Gromov--Witten theories; no Virasoro constraint for a
relative pair is an input.

\subsection{Geometry and coherent data}

\begin{notation}\label{sew:configurations}
All targets in a chain have complex dimension $n$.  Let $D$ be a smooth
projective variety, possibly disconnected, and let $N$ be a line bundle
on $D$.  We use the convention that $\PP(E)$ parametrizes lines in $E$.
Set
\[
 C=\PP_D(\cO_D\oplus N),\qquad
 D_l=\PP_D(N),\qquad D_r=\PP_D(\cO_D).
\]
The normals of $D_l,D_r$ in $C$ are $N^{-1},N$, respectively.  Thus a
chain, written in its geometric order, is
\begin{equation}\label{sew:chain}
 U\mid C_1\mid\cdots\mid C_m\mid V,
 \qquad C_i\cong C.
\end{equation}
Every displayed junction identifies a right section with a left section,
or identifies the corresponding divisor on an end.  Geometric order
will be kept distinct from the bipartite coloring used in
\cref{fp:local-algebra}.

At an internal junction, a \emph{surgery} separates the chain and
attaches a cap $Q$ to the part on the left and a cap $P$ to the part on
the right.  In both configurations below,
\[
 P=(C,D_r).
\]
The following table specifies the other data.  A ``short target'' is
the smooth fiber of the indicated capped segment; any number of necks
may be inserted into that segment.
\[
\begin{array}{c|c|c|ccc}
 &\text{main smooth target}&Q& U\mid Q&P\mid Q&P\mid V\\ \hline
 \text{ordinary}&X&(C,D_l)&U&C&V\\
 \text{blowup}&\Bl_S M&(Q,D)&M&Q&C
\end{array}
\]
In the ordinary row, $X$ is a smooth fiber of a given projective
degeneration with smooth total space and central fiber $U\cup_D V$;
$N=N_{D/U}$ and $N_{D/V}=N^{-1}$.  In the blowup row, $S\subset M$ is a
smooth center of codimension at least two,
\[
 U=\Bl_S M,\qquad D=\PP_S(N_{S/M}),\qquad
 N=\cO_D(-1),\qquad V=C.
\]
The absolute space underlying the right cap is
$Q=\PP_S(N_{S/M}\oplus\cO_S)$, with relative infinity divisor
$D=\PP_S(N_{S/M})$.  Here $D$ is also the exceptional divisor in $U$,
and $V$ is joined along $D_l$.  Empty centers require no blowup and may
be omitted.
\end{notation}

For a curve class $\gamma$ in a neck, write $\alpha=\pr_*\gamma$ and
$k_l=D_l\cdot\gamma$, $k_r=D_r\cdot\gamma$.  The divisor relation is
\begin{equation}\label{sew:flux}
 [D_r]-[D_l]=\pr^*c_1(N),\qquad
 k_r-k_l=\int_\alpha c_1(N).
\end{equation}
At joining sections the intersection numbers are the nonnegative total
contact orders of the relative maps.  At a nonjoining section the
intersection number is an ordinary divisor degree and may have either
sign.  For disconnected $D$, all formulas and, when required, their
degree conditions are imposed on its individual components.

\begin{lemma}[Realization and common local geometry]\label{sew:geometry}
For either configuration in \cref{sew:configurations}, chains of
arbitrarily large length and all unions of capped subchains obtained
from them have projective semistable realizations with smooth total
space.  They can be colored with two colors so that every junction has
oppositely colored ends.  On any fixed region disjoint from the
surgeries, the component and double-stratum diagrams, boundary normal
identifications, and first-page residue operations can be identified
across all the realizations.
\end{lemma}

\begin{proof}
In the ordinary configuration start with the given two-piece family.
After a finite base change, resolve the chain singularity; in a
transverse chart this is the projective semistable resolution of
$xy=t^a$.  It inserts the required copies of $C$.  The resolution is
obtained by blowups, so is projective, and its total space is smooth.
The same construction applied to deformation to the normal cone of the
divisor $D\subset U$, or $D\subset V$, gives the capped segments with
smooth fibers $U$ and $V$.  Deformation to the normal cone of either
section of $C$ gives the middle segments with smooth fiber $C$.

For the blowup configuration, deformation to the normal cone of
$D\subset U$ supplies the main chain, with smooth fiber $U$.
Deformation of $M$ to the normal cone of $S$ has central fiber
$U\cup_D Q$ and gives the left capped segment.  In $Q$ the normal of
its infinity divisor $D$ is $\cO_D(1)=N^{-1}$; deformation to the normal
cone of that divisor gives $P\mid Q$, with smooth fiber $Q$.  The
remaining segment is obtained by the divisor normal-cone degeneration
of $C$.  Inserting further levels by the same base change and
resolution gives any required length.  Take disjoint unions of these
families to realize several capped segments simultaneously.

Color the original chain alternately.  At a cut, give each new cap the
color opposite to its neighbor.  Each color is consequently a disjoint
union of smooth components, and the central fiber is a two-sided
degeneration without triple intersections.  This is the form to which
the two-sided degeneration formula applies.

Away from a surgery the smooth pairs and normal bundles have not
changed.  Their logarithmic charts factor the base generator at the
same double divisors.  Units in these local equations do not change
the graded residue maps.  The ordered configuration spaces in
\cref{ev:restriction} are constructed from these same divisors and
normal bundles.  Their projected incidence neighborhoods, and the
pullback, cup and trace operations on those neighborhoods, therefore
agree by \cref{fp:local-algebra}.  This assertion concerns the local
operations before passage to cohomology and makes no choice of global
representatives for surviving cohomology classes.
\end{proof}

Fix coherent admissible singles, in the sense of
\cref{sec:conventions}, and divisor degrees in these realizations.
The singles come from algebraic classes on the relevant total
families, allowing complex linear combinations, or from homogeneous
cohomology classes of fixed smooth projective sources through
algebraic correspondences extending over those total families and
inducing flat Hodge maps on their smooth loci.
Coherence means that their restrictions to every common component and
stratum agree, with the prescribed Hodge shifts.  For source-supplied
singles this agreement is required for the specified extending
algebraic correspondences with their cohomology slots exposed, before
applying the fixed inputs.  Retain these correspondences until
specialization and then apply their fixed inputs.  An end-supported
class is assigned zero on a capped segment not containing that end.
In particular, zero-section Gysin classes on a cap have fixed smooth
projective sources and are treated with their actual Gysin shifts.
For the target first Chern class use
$-c_1(\omega_{\mathcal Y/\Delta})$; on a component $Y$ it restricts to
$c_1(TY)-[\partial Y]$.  Thus the powers of the Chern-class operation in
the positive operators have coherent logarithmic restrictions on all
common pieces.

In the ordinary configuration, choose an integral relative
polarization on the original family.  Its degree, denoted by $q$, is
ample on $U$ and $V$ and is pulled back from $D$ on inserted necks and
caps.  We also denote its formal variable by $q$.  In the blowup
configuration, choose an ample integral class on $M$ and use its
pullback to $U$, its restriction to $S$, and the further pullbacks on
the bundle pieces.  Every effective contribution has nonnegative
counting degree.  This counting class need not be ample on a bundle
piece.

One may record additional coherent numerical divisors.  In the
ordinary configuration they extend over the original degeneration;
their common seam restrictions are pulled back to the necks and caps.
It suffices to choose lifts in the central restriction diagram, modulo
the usual opposite seam twists.  In the blowup configuration the
additional divisors used away from the final end are pulled back from
$M$ and $S$.  Give their variables invertible fugacities, so fixed
degree shifts in the cap coefficients are permitted.

There is one further degree condition in the blowup configuration:
retain the degree against the nonjoining section $D_r\subset V$,
component by component if necessary.  This section continues to the
exceptional divisor on the main smoothing.  Use an extension supported
at that final end, with zero restriction on all other components.
Consequently this degree is still measured on the segment containing
$V$ after any surgery; none of the inserted caps contributes to it.
The degrees against the nonjoining $D_l$ sections of inserted $P$ caps,
used in \cref{cap:identity}, are separate conditions.  They continue as
divisors on the corresponding short smooth targets, including a
segment of type $P\mid Q$.

Here the supported extensions can be constructed in the normal-cone
families themselves.  If $T\subset Y$ is the center, the strict
transform of $T\times\Delta$ in
$\Bl_{T\times\{0\}}(Y\times\Delta)$ is a constant family with central
fiber the zero section in its cap.  Gysin from this family extends
every class from the constant source $\HH(T)$, with specialization
supported on that zero section.  For a divisor center it also supplies
the nonjoining-section degree class.  When both ends are capped, the
centers on the short smoothing are disjoint: they are the two sections
of $C$ ordinarily, and the infinity divisor and zero section of $Q$ in
the blowup middle segment.  The two normal-cone constructions can be
performed simultaneously along these disjoint centers; further chain
insertions do not meet the zero sections.  Thus both sets of supported
extensions can be chosen independently.  Relations between their
classes on a smooth fiber do not prevent these choices of extensions
or the use of labelled multilinear insertions.

\subsection{The tested transport statement}

Let $s\geq 1$.  Write $Z_Y^{(r)}$ for independent replicas of the
disconnected descendant potential of a smooth target $Y$, with
independent descendant, degree, and genus variables; in particular the
genus variables are $\h_1,\ldots,\h_s$.  Fix one replica, numbered $1$,
and a positive integer $k$.  An admissible test $\mathsf T$ means a
finite tensor test of the class in \cref{sec:conventions}: it uses
coherent admissible singles, first-Hodge-index functions, cups,
coevaluations between arbitrary slots, classical integrations, and
ordinary GW amplitudes in the other replicas.  Its exposed arities,
descendant orders and numerical tensor operations are fixed.  Each
test is first decomposed into homogeneous terms.  Coefficient
extraction in the independent replicas is part of the test.

\begin{theorem}[Tested transport]\label{sew:transport}
Consider either configuration of \cref{sew:configurations}, with the
coherent admissible singles and degree data just specified.  Suppose that all
three short target types satisfy the absolute Virasoro constraints in
every genus, every curve class, and with all cohomology insertions.
Then, for every $k>0$, every finite number $s$ of independent replicas,
and every admissible test $\mathsf T$, all coefficients of
\begin{equation}\label{sew:tested-error}
 \mathsf T\left(
   (L_{k,X}^{(1)}Z_X^{(1)})
       \otimes\bigotimes_{r=2}^{s} Z_X^{(r)}
 \right)
\end{equation}
vanish, where $X$ is the main smooth target.  The coefficients are
selected by the coherent degree data, with the final exceptional
degrees retained in the blowup configuration.  The operator acts only
on replica $1$, before the test is applied.
\end{theorem}

The theorem is deliberately stated with the specified degree tests.
When these distinguish all relevant curve classes, it holds per curve
class; the curve-class verification in the application is made in
\cref{ind:degrees}.  The proof will not identify individual lifted
homology classes in different smoothings.

\subsection{Uniform preparations at the switches}

Fix the test in \eqref{sew:tested-error}, a desired coefficient in each
$\h_r$, and a finite counting-degree bound $d_r$ in each replica.  We
may use their maximum $d$ for all preparations.  Fix also the
additional degree coefficients, including the final exceptional end
degrees when present.  All calculations below are in the
coefficientwise completion of \cref{cap:identity}, modulo
$q_r^{d_r+1}$ in replica $r$.  A claim of equality always refers to
this completion, before the final coefficient is read.

\begin{lemma}[Preparations valid for every subset]\label{sew:preparation}
For any fixed finite ordered list of switches, one can choose a
contact bound and cap identity test at each switch so that the
following hold simultaneously for every subset of surgeries:
\begin{enumerate}[label=\textup{(\roman*)}]
 \item all incoming states lie in the chosen bounded domain;
 \item the chosen rows themselves bound outgoing contact, even when a
 positive-operator summand acts on their auxiliary labels;
 \item every coefficient calculation is finite in the degree and
 auxiliary-label orders under consideration and is bounded below in
 each genus variable;
 \item at a switch with no exposed operation, the total capped
 correspondence is the original contact-gluing correspondence on all
 outside states.
\end{enumerate}
Each replica has its own preparation and ordinary contact gluing.
\end{lemma}

\begin{proof}
This is the geometric application of \cref{cap:identity,cap:uniform};
we spell out the bounds that are used in the cancellation.  There are
constants $C_0,C_1$ such that, on effective classes,
\begin{equation}\label{sew:degree-bounds}
 \int_\alpha c_1(N)\leq C_0q(\alpha),\qquad
 D\cdot\beta_U\leq C_1q(\beta_U).
\end{equation}
Ordinarily these follow from ampleness on $D$ and $U$.  For the blowup
configuration, $-N$ is relatively ample over $S$, and minus the
exceptional divisor is relatively ample over $M$.  Adding sufficiently
large pullbacks of the chosen ample classes gives the inequalities.
At the initial relative end the contact is nonnegative, so the second
inequality bounds its total.  Through an unchanged neck,
\eqref{sew:flux} bounds any increase of contact by the counting degree
spent on that neck.

At a surgically inserted $P$, the preparation imposes one of a bounded
list of nonnegative initial fluxes $h$.  Its outgoing total is
\begin{equation}\label{sew:p-flux}
 a=h+\int_\alpha c_1(N)\leq h+C_0d.
\end{equation}
This conclusion uses only the imposed degree condition.  Altering a
descendant, applying a grading or a Chern-class power to an auxiliary
label, or removing labels to a classical operation does not change
that condition.  Choose the bound at the first switch from
\eqref{sew:degree-bounds}.  Choose the next one using the largest of
the unchanged-neck bound and all outgoing bounds in
\eqref{sew:p-flux} from the preceding preparation.  Continue from left
to right, taking the maximum over the finitely many earlier choices.
This proves uniformity over all subsets without requiring the
identity property at a busy switch.

For clarity, the zero-counting-degree exceptional classes have not
been removed in this argument.  If $\alpha\ne0$ is effective and
$q(\alpha)=0$ in the blowup configuration, it is vertical over $S$ and
$\int_\alpha c_1(N)<0$.  Such terms strictly lower the outgoing
contact.  They are included in the decreasing-contact part of the
inverse construction in \cref{cap:identity}.

In the final $V$ of the blowup configuration, let $e$ be its recorded
nonjoining section degree and let $a$ be the incoming contact.  Then
\[
 \int_\alpha c_1(N)=e-a.
\]
The fixed $e$ and bounded nonnegative $a$ give the lower bound missing
from a counting-degree bound alone.  Together with relative
ampleness, this bounds the effective horizontal degrees there.  The
imposed $h$ supplies the corresponding bound at inserted ends.  At a
$Q$ end, bounded incoming contact and bounded base degree bound the
fiber degree.  These are precisely the end conditions used in
\cref{cap:uniform} to bound the effective degrees on all intervening
pieces.  Extra divisor fugacities consequently have finite support in
the bounded blocks.

The cap inverse has nonnegative counting-degree powers and finitely
many selected auxiliary coefficient rows at these orders.  Its genus
coefficients are Laurent series bounded below.  The same holds for
the virtual sums with these bounds: the number of negative-genus-power
components is bounded at fixed degree and marking order, unmarked
constant genus-zero and genus-one components being unstable.  Higher
genus unmarked constants have positive genus power.  This proves the
needed coefficientwise finiteness, including for separate genus
variables.  Finally \cref{cap:identity} is an equality on the full
bounded relative state space, so it gives assertion (iv) before any
contraction with the outside data.
\end{proof}

The auxiliary rows in this lemma need not have the same arity.  Their
ordinary markings will be preintegrated in the evaluation
correspondences.  None of the forthcoming bounds on exposed positions
will depend on this variable arity.

\subsection{The alternating sum and its local terms}

We prove \cref{sew:transport}.  Choose a long chain and a set
$I=\{1,\ldots,b\}$ of widely separated internal switches.  The number
and separation will be fixed below using only the exposed test and
operator data, before selecting any cap rows.  For $A\subset I$, let
$X_A$ be the smooth disjoint union obtained by surgery at precisely
the switches in $A$; let $X_{\varnothing}=X$.  If $A\ne\varnothing$,
\begin{equation}\label{sew:short-unions}
 X_A\cong
 \begin{cases}
 U\amalg C^{\amalg(|A|-1)}\amalg V,&\text{ordinary},\\
 M\amalg Q^{\amalg(|A|-1)}\amalg C,&\text{blowup},
 \end{cases}
\end{equation}
where $Q$ in the second line denotes its absolute total space.
These identifications concern the smooth target types, with the
coherent extensions of labels and degrees specified above.

At every switch in $A$ and in every replica insert the prepared cap
tests of \cref{sew:preparation}.  A row choice $\rho$ specifies actual
ordinary cap descendants, their graded coefficients, and the required
degree projections.  Write $c_{A,\rho}$ for its external coefficient,
$\mathsf D_{\rho}^{(r)}$ for labelled extraction of its auxiliary
descendants in replica $r$, and $\mathsf P_{A,\rho}$ for its degree
projections and compensating shifts.  Let $\mathsf T_A$ be the fixed
original test interpreted on $X_A$ with the chosen coherent data.
Define
\begin{equation}\label{sew:subset-error}
 \mathcal E_A=
 \sum_{\rho} c_{A,\rho}\,\mathsf P_{A,\rho}\,
 \mathsf T_A\left(
  \mathsf D_{\rho}^{(1)}(L_{k,X_A}^{(1)}Z_{X_A}^{(1)})
  \otimes\bigotimes_{r=2}^{s}
       \mathsf D_{\rho}^{(r)}Z_{X_A}^{(r)}
 \right).
\end{equation}
All descendant extractions are evaluated with the remaining
auxiliary variables zero.  Any tensor-valued row coefficients in this
notation are contracted in their fixed graded order.

The order of operations in \eqref{sew:subset-error} is essential:
$L_k$ acts on the full descendant potential of replica $1$ before its
auxiliary descendants are extracted.  Thus it acts on auxiliary cap
labels just as on original labels, including a pair used by the
classical quadratic term.  It does not act on the external
$c_{A,\rho}$ or on any descendant variable in another replica.  In
particular no differentiation of the degree or genus series defining
the cap inverse is intended.  Multilinear extraction with distinct
label variables implements all marking multiplicities.

For $A\ne\varnothing$ one has
\begin{equation}\label{sew:nonempty-zero}
 \mathcal E_A=0.
\end{equation}
Indeed, for a disjoint union of targets, the potential is the product
of the potentials and the Virasoro operator is their sum, since the
cohomology and pairing are orthogonal direct sums.  By the hypotheses
and \eqref{sew:short-unions}, its full error therefore vanishes for all
descendant labels.  The auxiliary extractions, degree projections,
and graded tensor tests, even those supplied by other GW replicas,
preserve that zero identity.  No constraint is required for any of
the extra test replicas.  It remains to prove
\begin{equation}\label{sew:alternating-zero}
 \sum_{A\subset I}(-1)^{|A|}\mathcal E_A=0.
\end{equation}

\begin{lemma}[A bound on busy switches]\label{sew:busy-bound}
There are constants $B$ and $r_0$, depending on the fixed test, $k$,
and $n$, but independent of chain length, surgery subset, contact
bounds and auxiliary row arities, with the following property.
Each local summand in the expansion of any $\mathcal E_A$ has at most
$B$ busy switches if switches and ends are separated by more than
$2r_0+2$ necks.  All its cohomological operations occur in recorded
incidence neighborhoods of radius at most $r_0$.  A switch is busy if
its surgery region meets one of these neighborhoods or if an
auxiliary label at that switch is used nonordinarily by $L_k$.
\end{lemma}

\begin{proof}
Expand the fixed operator coefficient by \cref{conv:positive-rule}.
There are finitely many tensor operation patterns.  Record the fixed
original slots and every additional special GW slot.  Record also
any auxiliary slot moved to a classical operation or otherwise
modified by this summand of $L_k$.  There are at most two such
exceptional auxiliary slots: the linear part acts on one existing
slot, the classical quadratic part can use two, and the second-order
part creates its own pair of slots.  A modified auxiliary class may
be evaluated using its fixed, Gysin-shifted first Hodge index and the
coherent Chern-class cup operation.  Its support remains at its cap;
it is nevertheless counted as a position.

All other auxiliary labels are ordinary insertions from constant
cap sources.  Integrate them into the GW correspondence, retaining
only the exposed evaluations just listed.  The construction of
\cref{ev:restriction} applies with precisely that retained list;
\cref{ev:groups} permits the other group of labels to be preintegrated
without altering it.  The dimensions of the recorded products thus
have a bound determined by the original expression and the two
exceptional slots, independently of the number of auxiliary labels.

Now apply \cref{fp:local-algebra}.  Every first-page Casimir summand
has a single local vertex or edge support, even when its two legs
enter different replicas.  All pulls, cups and traces involve bounded
incidence neighborhoods of the recorded slots.  Finitely many
operation patterns on these bounded products give a uniform bound
on their number and on a radius $r_0$.  Increase the bound to include
the finitely many projected positions from all replicas and all
classical operations.  With the stated switch separation, each such
neighborhood can affect at most one switch; without this separation
one could instead multiply by a fixed bound for its number of
vertices.  Counting inserted-cap positions at their insertion switch
therefore gives a constant $B$ of the required kind.  The argument
does not count an ordinary preintegrated cap marking as a recorded
factor, so neither the number of rows nor their arities enters $B$.
\end{proof}

Choose $b>B$, the indicated separation, and enough necks at both ends.
Then make the preparations of \cref{sew:preparation}.  Expand every
summand of \eqref{sew:alternating-zero} by the two-sided component rule
of \cref{ev:restriction} and by the first-page calculation of
\cref{fp:local-algebra}.  The appropriate color classes may be
disconnected.  In particular no condition that distant recorded
positions lie on the same connected smooth target or on the same
connected domain is imposed on the individual local summands.

\subsection{Cancellation of grouped correspondences}

\begin{lemma}[Equality at an idle switch]\label{sew:idle-equality}
Fix a local operation pattern and a switch $j$ outside its busy set.
Keep its recorded component assignments and all data off the surgery
region at $j$.  Sum all relative-map histories, auxiliary row choices
and degree decompositions inside that region.  The resulting
correspondence is equal with $j$ uncut and with $j$ cut and capped.
This is an equality with the retained aligned evaluations and
ordinary cotangent classes, before the remaining cohomological
operations and tests are applied.
\end{lemma}

\begin{proof}
At an idle switch there is no recorded position on either inserted
cap, and no auxiliary label there is modified by $L_k$.  By
\cref{sew:geometry}, all retained local component and stratum data
agree on the two sides of the comparison.  The component restriction
formula of \cref{ev:restriction} is applied with these assignments
fixed.  Each color is a disjoint union of ends, necks, and caps, so
every connected component of a side source lies over one connected
component of that color.  In each color, collect all source components
mapping to the retained outside pieces into one group, with all of
that color's recorded positions.  Use separate groups for the cap
pieces at $j$.  The retained group is allowed to be disconnected.

Apply \cref{ev:groups} with this grouping.  Its common-refinement
comparison preserves the full joint evaluation of the retained group
into its aligned configuration space.  The retained ordinary marks
stabilize their own levels, so their cotangent classes are preserved
as well.  Projection formula integrates the ordinary cap groups,
which have no recorded positions.  The outside pairs and their
recorded data agree on the two sides; hence this leaves the same
outside relative correspondence, with two uncontracted contact-state
inputs at $j$ in each replica.  The only remaining difference is the
bilinear kernel contracting these states.  With $j$ uncut it is the
ordinary contact kernel $G_{\h_r}$ of \cref{cap:identity}.  With $j$
cut it is the sum of the
two cap kernels with their prepared coefficients and degree
conditions.  By \cref{cap:identity}, these kernels are equal on every
allowed incoming state and on every compatible outgoing state.
\Cref{sew:preparation} ensures that those domains include all terms
under consideration, including histories in which other switches
are busy.

The comparison is made separately in each replica, with its own
$q_r$, $\h_r$, states and contact joins.  Tensor the resulting
equalities in the fixed order.  Equality on the full relative state
spaces permits subsequent contraction with arbitrary outside
tensors, including tensors correlating different replicas.  It
never requires gluing a contact from one replica to a contact in
another.  All contact-order factors, finite permutation quotients,
and one genus factor per joined root are already present in
$G_{\h_r}$.  Consequently changes in domain connectedness or in the
distribution of handles require no supplementary sign or condition.
The only signs for odd labels are the categorical tensor signs fixed
throughout.

Finally the histories are summed using the extended degree tests,
with the compensating shifts of the cap coefficients.  Their
coefficientwise sums exist by \cref{sew:preparation}.  One has not
chosen individual homology classes to match across different
smoothings.  The equality is therefore one of the full weighted
virtual correspondences needed as input to the remaining local
operations, as asserted.
\end{proof}

We finish the proof of \cref{sew:transport}.  Group first by the
position pattern, operator roles, component assignments for each
recorded evaluation, and incidence and orientation data in the
retained neighborhoods.  These data determine the busy set.  Every
such pattern has an idle switch; let $j$ be its first idle switch in
the fixed left-to-right order.  Refine this grouping by the calculation
outside the region at $j$, and sum all relative-map histories,
ordinary auxiliary labels, and inverse-matrix row choices inside
that region.  These internal choices are not part of the recorded
position pattern.  The sums are legitimate coefficientwise by
\cref{sew:preparation}; the resolved products on the two sides have
the same recorded factors even if the numbers of preintegrated cap
markings differ.

Replacing $A$ by $A\mathbin{\triangle}\{j\}$ changes only the unrecorded
history at $j$.  It preserves the position pattern and all operator
roles; in particular any exceptional auxiliary modifications at other
switches remain the same.  It consequently preserves the busy set
and its first idle switch.  This gives an involution on the grouped
comparisons.  By \cref{sew:idle-equality}, the two groups supply
identical tensors to every retained cohomological operation.  They
have opposite inclusion--exclusion signs, since the subset cardinality
changes by one.  Fixing the order of chain components, marked slots
and cap coefficient spaces once and for all also identifies the
graded permutation conventions on both sides.

This cancellation is between summed correspondences.  It does not
assert a bijection between individual stable maps or individual
inverse-matrix rows.  Empty boundary profiles, components without
ordinary insertions, and the classical term carrying two removed
labels are included: the first two are part of the disconnected cap
identity, and the last is part of the recorded local operation
pattern.  Ordinary marks are distinguished in multilinear
coefficient extraction, so changing a row's number of preintegrated
marks introduces no unaccounted placement multiplicity.

We have proved \eqref{sew:alternating-zero}.  All its nonempty-subset
terms vanish by \eqref{sew:nonempty-zero}, leaving
$\mathcal E_{\varnothing}=0$.  This is exactly the prescribed
coefficient of \eqref{sew:tested-error}.  The degree truncations,
genus coefficients, additional degree coefficients and test were
arbitrary.  The tested transport theorem follows.
